% Aula 01 --- o preditor ótimo e o erro irredutível.
% Incluído por "Lista de revisao 01-06.tex" e, se sorteado, pela prova.
No início do curso, vimos uma afirmação forte: de todas as funções que você poderia usar
para prever $Y$ a partir de $\X$, a melhor é a função de regressão
$r(\x)=\E[Y\mid\X=\x]$. Não a melhor entre as lineares, nem a melhor entre as que
o seu método consegue produzir --- a melhor, ponto.

Uma afirmação dessas pede demonstração, e ela é mais curta do que parece. Todo o
trabalho está num truque que você vai reencontrar várias vezes no curso: somar e
subtrair a quantidade certa dentro de um quadrado.
\begin{itens}
  \item \pts{1,2} Seja $g:\R^p\to\R$ uma função qualquer com
        $\E[g(\X)^2]<\infty$. Some e subtraia $r(\X)$ dentro do quadrado e mostre
        que
        \[
          \E\big[(Y-g(\X))^2\big]
            = \E\big[(Y-r(\X))^2\big] + \E\big[(r(\X)-g(\X))^2\big].
        \]
        O trabalho está em mostrar que o termo cruzado desaparece; condicione em
        $\X$ e pergunte-se o que $\E[Y-r(\X)\mid\X]$ vale.
  \item \pts{0,9} Olhe para a identidade que você acabou de provar e responda: por
        que ela mostra que $r$ é o minimizador? Em seguida, mostre que o valor
        mínimo é $\risco(r)=\E[\Var(Y\mid\X)]$.
  \item \pts{0,8} Esse mínimo tem nome: \textbf{erro irredutível}. O nome sugere
        que não há como baixá-lo, e é verdade --- mas a razão não é que os métodos
        atuais sejam ruins. Explique de onde vem essa impossibilidade, e diga o que
        teria de mudar \emph{no problema} para o número diminuir.
  \item \pts{1,1} Na Lista Teórica~01, mostramos que o melhor preditor
        \emph{constante} tem risco $\Var(Y)$. Use a identidade de (a) para escrever
        a diferença $\Var(Y)-\risco(r)$ como uma esperança, e interprete: que
        pergunta prática essa quantidade responde?
\end{itens}

\begin{solucao}
Suponha, ao longo de todo o exercício, que $\E[Y^2]<\infty$: é o que garante que
todas as esperanças abaixo existam. Em particular $r(\X)$ também tem segundo momento
finito, porque $\E[r(\X)^2]=\E\big[(\E[Y\mid\X])^2\big]\le\E\big[\E[Y^2\mid\X]\big]=\E[Y^2]$
(desigualdade de Jensen, aplicada condicionalmente a $\X$).

\smallskip
\textbf{(a)} Some e subtraia $r(\X)$ e expanda o quadrado:
\[
  (Y-g(\X))^2 = \big(Y-r(\X)\big)^2
     + 2\big(Y-r(\X)\big)\big(r(\X)-g(\X)\big)
     + \big(r(\X)-g(\X)\big)^2 .
\]
As três parcelas são integráveis: as duas quadráticas porque $Y$, $r(\X)$ e $g(\X)$
têm segundo momento finito, e o termo cruzado pela desigualdade de Cauchy--Schwarz,
$\E\abs{UV}\le\sqrt{\E[U^2]\,\E[V^2]}$. Podemos, então, tomar esperança termo a
termo, e tudo se resume a mostrar que o termo cruzado tem esperança zero.

Condicione em $\X$ e use a lei da esperança total:
\[
\begin{aligned}
  \E\big[(Y-r(\X))(r(\X)-g(\X))\big]
  &= \E\Big[\,\E\big[(Y-r(\X))(r(\X)-g(\X))\,\big|\,\X\big]\Big]\\
  &= \E\Big[\big(r(\X)-g(\X)\big)\,
       \underbrace{\E\big[Y-r(\X)\,\big|\,\X\big]}_{=\;r(\X)-r(\X)\;=\;0}\Big]
  = 0 .
\end{aligned}
\]
O passo do meio é a regra "o que é função de $\X$ sai da esperança condicional em
$\X$": $r(\X)-g(\X)$ depende só de $\X$, e condicionar em $\X$ é tratá-lo como
conhecido. O zero vem da \emph{definição} de $r$: $\E[Y\mid\X]=r(\X)$.

Em palavras: o erro $Y-r(\X)$ que sobra depois de acertar a média condicional é
\textbf{ortogonal a toda função de $\X$} --- nenhuma informação contida em $\X$ ajuda a
prevê-lo. Guarde a frase: o Exercício~\ref{ex:02-oraculo} tem uma versão dela para
as funções lineares.

\smallskip
\textbf{(b)} A identidade de (a) diz que
\[
  \risco(g) \;=\; \underbrace{\risco(r)}_{\text{não depende de } g}
   \;+\; \underbrace{\E\big[(r(\X)-g(\X))^2\big]}_{\ge\ 0} .
\]
Logo $\risco(g)\ge\risco(r)$ para toda $g$: $r$ é minimizadora. E é a única, a menos
de probabilidade zero: a igualdade vale se e só se $\E[(r(\X)-g(\X))^2]=0$, isto é,
se $g(\X)=r(\X)$ com probabilidade $1$. (Alterar $g$ num conjunto de valores de $\x$
que nunca ocorrem não muda risco nenhum.)

O valor mínimo sai de novo pela esperança total. Dado $\X$, $Y-r(\X)$ é o desvio de
$Y$ em torno da sua própria média condicional, e a esperança condicional do seu
quadrado é, por definição, a variância condicional:
\[
  \risco(r)=\E\big[(Y-\E[Y\mid\X])^2\big]
   =\E\Big[\,\E\big[(Y-\E[Y\mid\X])^2\;\big|\;\X\big]\Big]
   =\E\big[\Var(Y\mid\X)\big].
\]

\smallskip
\textbf{(c)} A impossibilidade é matemática, e não tecnológica. O item (b) mostrou que
\textbf{nenhuma} função de $\X$ tem risco menor que $\E[\Var(Y\mid\X)]$ --- nem mesmo
$r$, conhecida exatamente, com $n=\infty$ e erro de estimação zero. E essa quantidade
não menciona método nenhum: é uma propriedade da distribuição conjunta de $(\X,Y)$.
Ela mede quanto $Y$ ainda varia entre indivíduos com \emph{o mesmo} $\x$. Dois
pacientes com idade, peso e exames idênticos, mas desfechos diferentes, recebem
forçosamente a mesma previsão, porque o preditor recebeu a mesma entrada nos dois
casos.

Para baixar o número é preciso mudar o problema: \textbf{medir mais}. Se $Z$ é uma
covariável nova, a lei da variância total, aplicada condicionalmente a $\X$, dá
\[
  \E\big[\Var(Y\mid\X)\big]
  = \E\big[\Var(Y\mid\X,Z)\big]
    + \E\Big[\Var\big(\E[Y\mid\X,Z]\;\big|\;\X\big)\Big]
  \;\ge\; \E\big[\Var(Y\mid\X,Z)\big],
\]
com igualdade só quando $Z$ não muda a média condicional, $\E[Y\mid\X,Z]=\E[Y\mid\X]$.
Acrescentar informação nunca aumenta o erro irredutível, e o diminui sempre que a
informação for relevante. Isso vale na população; com $n$ finito, cada covariável
nova também custa variância de estimação (Aula~02), de modo que medir mais é
necessário, mas não suficiente, para prever melhor.

\smallskip
\textbf{(d)} Aplique a identidade de (a) com $g\equiv\E[Y]$. O lado esquerdo é o risco
do melhor preditor constante, $\Var(Y)$. E como $\E[r(\X)]=\E\big[\E[Y\mid\X]\big]=\E[Y]$,
pela esperança total, a última parcela é uma variância:
\[
  \Var(Y) = \risco(r) + \E\Big[\big(r(\X)-\E[r(\X)]\big)^2\Big]
  \qquad\Longrightarrow\qquad
  \boxed{\;\Var(Y)-\risco(r)=\Var\big(r(\X)\big)\;}
\]
É a lei da variância total, $\Var(Y)=\E[\Var(Y\mid\X)]+\Var(\E[Y\mid\X])$, obtida por
outro caminho.

A pergunta que isso responde é: \textbf{vale a pena usar as covariáveis?} O lado
esquerdo é o máximo que se pode ganhar trocando o preditor constante pelo melhor
preditor possível; o direito diz que esse ganho é exatamente o quanto a resposta
\emph{esperada} muda de um perfil de covariáveis para outro. Dividido por $\Var(Y)$,
vira a fração da variância de $Y$ que as covariáveis conseguem explicar no melhor dos
casos: um $R^2$ populacional, teto para o $R^2$ de qualquer modelo.

Os dois extremos ajudam a fixar. Se $\Var(r(\X))=0$, a função de regressão é
constante: as covariáveis não carregam informação sobre a média de $Y$, e prever pela
média é o melhor que existe --- não há modelo a construir. Se $\risco(r)=0$, $Y$ é
função de $\X$, e toda a sua variância é explicável.
\end{solucao}
